\section{章节罗盘：战争如何改写筹饷与权力的路径}
\label{sec:nineteenth-conflicts-compass}

条约冲击没有取代清廷原有的危机，反而与内战、河患、难民和边务叠加。本章将太平、捻军以及随后西北、西南的战争放在粮道、税源、港口和军械的共同网络中考察：一座城的失守、一条河道的改向或一次借饷，都会改变远处军队与地方社会的选择。

本章的中心问题也不是“中央弱、地方强”的简单替代。湘军、淮军、厘金、捐输和省级中介让清廷获得继续作战和重建的资源，同时改变了兵权、信息和财源的实际归属。收复与新式机构并非战争之外的插曲；它们是危机财政的产物。结尾将把这种重组带到改革时代：国家必须借助地方网络，才能尝试重新集中资源。
